\documentclass[12pt]{article}
\usepackage{amsmath}
\usepackage{enumitem}
\usepackage{fancyhdr}
\usepackage[margin=1in]{geometry}
\usepackage{courier}
\usepackage{listings}
\usepackage{xcolor}

\pagestyle{fancy}
\fancyhf{}
\rhead{CPSC 250 -- Python for Data Manipulation -- Summer 2026 Test 1}
\lhead{Name:\underline{\hspace{6cm}}}
\rfoot{Page \thepage}

\lstset{
  basicstyle=\ttfamily\small,
  breaklines=true,
  keywordstyle=\color{blue},
  commentstyle=\color{gray},
  stringstyle=\color{orange},
  showstringspaces=false,
}

\begin{document}

\begin{center}
    \Large \textbf{CPSC 250 -- Python for Data Manipulation -- Summer 2026 Test 1} \\
    \normalsize Topics: Primitive Data Types, Lists \& Tuples, Branching, Loops
\end{center}

\vspace{0.5cm}

\section*{Part I -- Multiple Choice (10 points)}

Circle the best answer. Each question is worth 2 points.

\begin{enumerate}[label=\arabic*.]
    \item What is the result of \texttt{17 \% 5}?
    \begin{enumerate}[label=\Alph*.]
        \item 3
        \item 2
        \item 5
        \item 17
    \end{enumerate}

    \item Which of the following creates a tuple?
    \begin{enumerate}[label=\Alph*.]
        \item \texttt{\char`\[4, 5, 6\char`\]}
        \item \texttt{\{4, 5, 6\}}
        \item \texttt{(4, 5, 6)}
        \item \texttt{"4, 5, 6"}
    \end{enumerate}

    \item What is the result of \texttt{float("7") + 2}?
    \begin{enumerate}[label=\Alph*.]
        \item \texttt{9}
        \item \texttt{9.0}
        \item \texttt{"72"}
        \item Error
    \end{enumerate}

    \item Which condition checks whether \texttt{age} is less than 18 or greater than 65?
    \begin{enumerate}[label=\Alph*.]
        \item \texttt{age < 18 and age > 65}
        \item \texttt{age < 18 or age > 65}
        \item \texttt{18 <= age <= 65}
        \item \texttt{not age}
    \end{enumerate}

\newpage

    \item What is printed by the following code?
\begin{lstlisting}
data = [10, 20, 30]
print(data[-1])
\end{lstlisting}
    \begin{enumerate}[label=\Alph*.]
        \item \texttt{10}
        \item \texttt{20}
        \item \texttt{30}
        \item Error
    \end{enumerate}
\end{enumerate}

\section*{Part II -- Find the Error (15 points)}

Each of the following code snippets contains one or more errors. Identify and correct them.

\begin{enumerate}[resume,label=\arabic*.]
    \item
\begin{lstlisting}
scores = [85, 90, 95]
print(scores(1))
\end{lstlisting}

    \vspace{2.5cm}

    \item
\begin{lstlisting}
if x > 10
    print("Large")
\end{lstlisting}

    \vspace{2.5cm}

    \item
\begin{lstlisting}
total = 0
for i in range(5):
total += i
print(total)
\end{lstlisting}

    \vspace{2.5cm}
\end{enumerate}

\newpage

\section*{Part III -- Code Writing (20 points)}

Write Python code for each of the following tasks.

\begin{enumerate}[resume,label=\arabic*.]
    \item Write Python code that asks the user to enter 8 integers. Count how many of the numbers are even, then print the count.

    \vspace{11cm}

\newpage

    \item Write Python code that uses a loop to print all multiples of 5 from 5 through 50 inclusive, each on its own line.

    \vspace{8cm}
\end{enumerate}

\newpage

\section*{Part IV -- Code Comprehension and Commenting (15 points)}

The following program is complete but has no comments. Write comments to explain what each line does.

\begin{lstlisting}
# ________________________________________

numbers = [4, 8, 1, 6, 9]

# ________________________________________

largest = numbers[0]

# ________________________________________

for n in numbers:

# ________________________________________

    if n > largest:
        largest = n

# ________________________________________

print("Largest value:", largest)
\end{lstlisting}

\vspace{3cm}

\newpage

\section*{Part V -- Short Answer (10 points)}

\begin{enumerate}[resume,label=\arabic*.]
    \item Explain the difference between \texttt{=} and \texttt{==} in Python. Give an example of each.

    \vspace{3cm}

    \item What is the difference between a \texttt{for} loop and a \texttt{while} loop? Describe a situation where each would be appropriate.

    \vspace{3cm}
\end{enumerate}

\section*{Total: \underline{\hspace{2cm}} / 70}

\end{document}
